\section{EXTENSION OF A HARDY FIELD}

Given a Hardy field \(\mathfrak{K}\) we shall see how one can form new Hardy fields \(\mathfrak{K}' \supset \mathfrak{K}\) such that \(\mathfrak{K}' / \mathrm{R}_{\infty}\) is obtained by adjoining (in the algebraic sense of the term, cf.~Alg., V, §2) new elements, of a form which we shall make precise, to \(\mathfrak{K} / \mathrm{R}_{\infty}\) .

Lemma 1. Let \(a(x), b(x)\) be two continuous real functions that do not change sign on an interval \([x_0, +\infty[\) . If, on this interval, the function \(y(x)\) is continuous and differentiable, and satisfies the identity

\[
y ^ {\prime} = a y + b\tag{1}
\]

then there exists an interval \([x_{1}, +\infty[\) on which y does not change sign.

Indeed, let us put \(z(x) = y(x) \exp \left(-\int_{x_0}^x a(t)dt\right)\) (cf.~IV, p.~183); then, by (1), \(z'(x) = b(x) \exp \left(-\int_{x_0}^x a(t)dt\right)\) . If \(b(x) \geqslant 0\) for \(x \geqslant x_0\) then \(z\) is increasing on this interval, so either is \(< 0\) on all this interval, or is zero on an interval \([x_1, +\infty[\) , or else is \(> 0\) on an interval \([x_1, +\infty[\) ; since \(y\) has the same sign as \(z\) the proposition is proved in this case. The argument is similar if \(b(x) \leqslant 0\) for \(x \geqslant x_0\) .

Remark. This very elementary property does not extend to linear differential equations of order \textgreater{} 1; for example, the function \(y = \sin x\) satisfies \(y'' + y = 0\) , but changes sign on every neighbourhood of \(+\infty\) .

Lemma 2. Let \(a(x)\) and \(b(x)\) be two functions belonging to a given Hardy field \(\mathfrak{K}\) and \(y(x)\) a function satisfying the identity (1) on an interval \([x_0, +\infty[\) where \(a\) and \(b\) are defined and continuous. If \(p(u)\) is a polynomial in \(u\) whose coefficients are functions of \(x\) belonging to \(\mathfrak{K}\) , defined and differentiable on \([x_0, +\infty[\) , then there exists an interval \([x_1, +\infty[\) on which the function \(p(y)\) does not change sign.

The proposition is trivial if \(p(u)\) has coefficients which are identically zero on \([x_{0}, +\infty[\) , or if \(p(u)\) is of degree 0 in u, since any function in K is of constant sign on an interval \([x_{1}, +\infty[\) . Suppose that \(p(u)\) is of degree n \textgreater{} 0; the leading coefficient c of \(p(u)\) is then \(\neq 0\) on an interval \([\alpha, +\infty[\) ; one can thus write \(p(u) = c(u^{n} + c_{1}u^{n-1} + \cdots + c_{n})\) where \(c, c_{1}, c_{2}, \ldots, c_{n}\) are functions belonging to K and differentiable on \([\alpha, +\infty[\) ; so it is enough to prove the lemma for c = 1. We argue by induction on n; one has

\[
\begin{array}{c} \frac {d}{d x} (p (y)) = (a y + b) \bigl (n y ^ {n - 1} + (n - 1) c _ {1} y ^ {n - 2} + \dots + c _ {n - 1} \bigr) \\ + c _ {1} ^ {\prime} y ^ {n - 1} + \dots + c _ {n} ^ {\prime} = n a p (y) + q (y) \end{array}
\]

where \(q(u)\) is a polynomial of degree \(\leqslant n - 1\) , with coefficients in \(\mathfrak{K}\) . By hypothesis the functions \(na(x)\) and \(q(y(x))\) do not change sign on an interval \([\beta, +\infty[\) ; the lemma is thus a consequence of lemma 1.

THEOREM 1. Let \(a(x)\) and \(b(x)\) be two functions belonging to a given Hardy field \(\mathfrak{K}\) and \(y(x)\) a function satisfying (1) on an interval \([x_0, +\infty[\) . When \(r(u) = p(u) / q(u)\) runs through the set of rational fractions in \(u\) with coefficients in \(\mathfrak{K}\) such that \(q(y)\) is not identically zero on a neighbourhood of \(+\infty\) the set \(\mathfrak{K}(y)\) of functions \(r(y)\) forms a Hardy field.

Indeed, by lemma 2, there exists an interval \([x_1, +\infty[\) on which \(r(y)\) is defined, continuous, and does not change sign, from which it follows immediately that \(\mathfrak{K}(y)/\mathrm{R}_{\infty}\) is a field; moreover, since

\[
\frac {d}{d x} (r (y)) = r ^ {\prime} (y) y ^ {\prime} = r ^ {\prime} (y) (a y + b)
\]

(where \(r'(y) = \left(p'(y)q(y) - p(y)q'(y)\right)/\left(q(y)\right)^2\) is defined by hypothesis on a neighbourhood of \(+\infty\) ), the derivative of every function in \(\mathfrak{K}(y)\) belongs to \(\mathfrak{K}(y)\) , which proves that \(\mathfrak{K}(y)\) satisfies the conditions of def. 1 of V, p.~247.

It is clear that \(\mathfrak{K}(y) / \mathbb{R}_{\infty}\) is obtained by the algebraic adjunction to \(\mathfrak{K} / \mathbb{R}_{\infty}\) of the class of \(y\) modulo \(\mathbb{R}_{\infty}\) . Again one says that \(\mathfrak{K}(y)\) is obtained by adjoining \(y\) to \(\mathfrak{K}\) .

COROLLARY 1. If \(y\) is a function in \(\mathfrak{K}\) , not identically zero on a neighbourhood of \(+\infty\) , then \(\mathfrak{K}(\log |y|)\) is a Hardy field.

Indeed, \((\log |y|)' = y'/y\) is equal to a function in \(\mathcal{K}\) on an interval \([x_0, +\infty[\) .

COROLLARY 2. If \(y\) is any function in \(\mathfrak{K}\) , then \(\mathfrak{K}(e^y)\) is a Hardy field.

Indeed, \(\left(e^{y}\right)^{\prime} = e^{y}y^{\prime}\) , and \(y^\prime\) is equal to a function in \(\mathcal{K}\) on an interval \([x_0, +\infty[\) .

COROLLARY 3. If \(\mathfrak{K}\) contains the real constants, and if \(y\) is a function in \(\mathfrak{K}\) , not identically zero on a neighbourhood of \(+\infty\) , then \(\mathfrak{K}(|y|^{\alpha})\) is a Hardy field for every real number \(\alpha\) .

Indeed, \(\frac{d}{dx}\big(|y|^{\alpha}\big) = |y|^{\alpha}(\alpha y'/y)\) , and \(\alpha y'/y\) is equal to a function in \(\mathfrak{K}\) on an interval \([x_0, +\infty[\) .

Finally, we note that if \(y\) is the primitive of any function in \(\mathfrak{K}\) then \(\mathfrak{K}(y)\) is again a Hardy field.

